\chapter{Installation and daily use}
\label{ch:installation}

\section{Prerequisites}

ATTPCROOT is built against an existing installation of FairSoft and FairRoot. FairSoft provides ROOT, Geant4,
the virtual Monte Carlo (VMC) and their dependencies; FairRoot provides the task and run framework. The top-level
\code{CMakeLists.txt} requires FairRoot, ROOT, VMC, FairLogger, Geant4, Geant4VMC, Boost and HDF5, and looks for the
following optional packages:

\begin{center}
\small
\begin{tabular}{p{0.18\textwidth}p{0.74\textwidth}}
\toprule
Package & Used for \\
\midrule
\code{GENFIT2} & the GENFIT Kalman fitter (\code{AtGenfitter}, Chapter~\ref{ch:fitting}) \\
\code{CATIMA} & energy loss and multiple scattering (\code{AtELossCATIMA}); downloaded and built locally if it is
  not found \\
\code{Eigen3} & the unscented Kalman filter (OpenKF) \\
\code{CUDA}, \code{HiRAEVT} & optional components \\
\bottomrule
\end{tabular}
\end{center}

\section{First-time configuration}

The configure step needs to know where FairRoot and FairSoft are installed. Set two variables in the shell used
for the first configuration only:
\begin{verbatim}
export FAIRROOTPATH=/path/to/fairroot/install
export SIMPATH=/path/to/fairsoft/install
\end{verbatim}
Then configure out of source from the repository root, adding the HDF5 prefix if HDF5 is not part of FairSoft:
\begin{verbatim}
cmake -S . -B build -DCMAKE_PREFIX_PATH=/path/to/hdf5
cmake --build build -j10
\end{verbatim}
A successful configuration writes \code{build/config.sh}, which records the whole run-time environment. The two
variables above are captured in it and are not needed again.

\section{Every session}

Source the environment script from the repository root before building or running anything:
\begin{verbatim}
source build/config.sh
\end{verbatim}
It sets \code{ROOTSYS}, \code{VMCWORKDIR} (the repository root, used by macros to find \code{geometry/},
\code{parameters/}, \code{scripts/} and \code{resources/}), \code{ROOT\_INCLUDE\_PATH}, \code{LD\_LIBRARY\_PATH} and
the Geant4 data paths. Macros that use GENFIT also need its library path; the a1975 production scripts set it
explicitly:
\begin{verbatim}
export GENFIT=/path/to/GenFit
export LD_LIBRARY_PATH=$GENFIT/lib:$LD_LIBRARY_PATH
export ROOT_INCLUDE_PATH=$VMCWORKDIR/build/include:$ROOT_INCLUDE_PATH
\end{verbatim}

\paragraph{Building.} \code{cmake --build build -j10} rebuilds everything; \code{--target AtTools} rebuilds one
library. Public headers are copied to \code{build/include}.

\paragraph{Running a macro.} Macros are interpreted by ROOT's Cling. For production and remote work run them in
batch mode, with no graphics window, and with standard input closed so that ROOT cannot wait for input:
\begin{verbatim}
root -l -b -q 'unpackReco_multifit.C("run_0016")' < /dev/null > reco.log 2>&1
\end{verbatim}
Macros load the libraries they need with \code{gSystem->Load("libAtReconstruction.so")} and similar calls. If a
macro cannot resolve a framework symbol, suspect the environment or a missing library load before adding project
headers to the macro.

\paragraph{Running in parallel.} Long productions are split by run or by event range and run as background jobs. A
job started with a plain \code{\&} from an interactive shell can be killed when that shell exits; start long jobs
from a script with \code{nohup}. Each parallel simulation job needs its own random seed (Chapter~\ref{ch:simulation}).

\section{Tests}

Unit tests are built by default and placed in \code{build/tests/}:
\begin{verbatim}
cd build && ctest            # all tests
./build/tests/AtToolsTests   # one test binary
\end{verbatim}
At the time of writing three \code{AtPropagatorTest} tests fail; they use an external SRIM energy-loss table and are
discussed in Chapter~\ref{ch:tools}. Integration tests are ROOT macros in \code{macro/tests/}; they need geometry
files generated by \code{macro/tests/generateGeometry.sh} and, for some detectors, data that are not in the
repository. Chapter~\ref{ch:development} explains how to write tests.
